\diarytitle{0075}{3次対称行列の固有値・固有ベクトルと直交対角化}

$A = I + J$（$J$ はすべての成分が $1$ の $3\times3$ 行列）と書ける。$J$ の固有値は、階数 $1$・trace $3$ より $3, 0, 0$（$\lambda=3$ の固有ベクトルは $(1,1,1)^{\mathrm{T}}$、$\lambda=0$ の固有空間は $x+y+z=0$ を満たす平面）であるから、$A = I+J$ の固有値は
\[
\lambda = 1+3 = 4 \ (\text{単根}), \qquad \lambda = 1+0 = 1 \ (\text{重根})
\]
実際に検算すると、$\det(A - \lambda I)$ を直接展開しても
\[
\det(A-\lambda I) = -(\lambda-4)(\lambda-1)^{2}
\]
となり一致する。

$\lambda = 4$: 固有空間は $(1,1,1)^{\mathrm{T}}$ 方向。正規化して $\bm{p}_1 = \dfrac{1}{\sqrt{3}}(1,1,1)^{\mathrm{T}}$。

$\lambda = 1$: 固有空間は平面 $x+y+z=0$。この中で正規直交基底を（グラム・シュミットで）選ぶと、例えば
\[
\bm{p}_2 = \frac{1}{\sqrt{2}}(1,-1,0)^{\mathrm{T}}, \qquad
\bm{p}_3 = \frac{1}{\sqrt{6}}(1,1,-2)^{\mathrm{T}}
\]
（$\bm{p}_2 \cdot \bm{p}_3 = \frac{1}{\sqrt{12}}(1-1+0)=0$、$\bm{p}_2, \bm{p}_3$ とも $(1,1,1)$ に直交することも確認できる）。

したがって直交行列
\[
P = \begin{pmatrix}
\dfrac{1}{\sqrt3} & \dfrac{1}{\sqrt2} & \dfrac{1}{\sqrt6} \\[6pt]
\dfrac{1}{\sqrt3} & -\dfrac{1}{\sqrt2} & \dfrac{1}{\sqrt6} \\[6pt]
\dfrac{1}{\sqrt3} & 0 & -\dfrac{2}{\sqrt6}
\end{pmatrix}
\]
により
\[
\boxed{\; P^{\mathrm{T}} A P = \begin{pmatrix} 4 & 0 & 0 \\ 0 & 1 & 0 \\ 0 & 0 & 1 \end{pmatrix} \;}
\]

検算: $A\bm{p}_1 = (I+J)\bm{p}_1 = \bm{p}_1 + \sqrt3(1,1,1)^{\mathrm{T}} = \bm{p}_1 + 3\bm{p}_1 = 4\bm{p}_1$。また $J\bm{p}_2 = \bm{0}$（各成分の和が $0$ のため）より $A\bm{p}_2 = \bm{p}_2 + \bm{0} = 1\cdot\bm{p}_2$、同様に $A\bm{p}_3 = \bm{p}_3$。いずれも一致する。
